\section*{Capítulo \ref{ch:b2:elements-in-the-earth} --- Os elementos na Terra}

\begin{solution}{exo:b2:elements-in-the-earth:1}
O potássio é $2\times39.10/(2\times39.10+16.00) = 78.20/94.20 = 0.8301$ da massa do
\ce{K2O}, portanto $\num{5130}\times0.8301 = \qty{4259}{ppm}$ de K. O titânio é
$47.87/79.87 = 0.5993$ do \ce{TiO2}: $\num{27310}\times0.5993 = \qty{16370}{ppm}$ de Ti.
\end{solution}

\begin{solution}{exo:b2:elements-in-the-earth:2}
O alumínio é litófilo e refratário: a Terra silicática conservou-o todo, em proporções quase
condríticas (1.16). O sódio é litófilo, mas volátil (\qty{958}{K}): foi
perdido antes ou durante a acreção, restando 0.22. O níquel se condensa com o magnésio
(\qty{1353}{K}), portanto a volatilidade não explica 0.077: ele é siderófilo, e a maior parte dele
foi para o núcleo. O irídio é refratário (\qty{1603}{K}) e fortemente siderófilo:
0.0029, quase todo ele no núcleo.
\end{solution}

\begin{solution}{exo:b2:elements-in-the-earth:3}
$\mathrm{La}_N = 31/0.232 = 133.6$, $\mathrm{Yb}_N = 2.0/0.163 = 12.27$, e
$(\mathrm{La}/\mathrm{Yb})_N = 10.9$. A crosta é enriquecida em todas as terras raras, e
onze vezes mais nas leves que nas pesadas: ela é feita de fundidos, que
recebem com mais facilidade as terras raras leves incompatíveis.
\end{solution}

\begin{solution}{exo:b2:elements-in-the-earth:4}
$\mathrm{Sm}_N = 4.7/0.145 = 32.4$, $\mathrm{Eu}_N = 1.0/0.0546 = 18.3$,
$\mathrm{Gd}_N = 4.0/0.198 = 20.2$; $\mathrm{Eu}^* = \sqrt{32.4\times20.2} = 25.6$ e
$\mathrm{Eu}/\mathrm{Eu}^* = 18.3/25.6 = 0.72$, uma anomalia negativa. O plagioclásio, que
aceita o \ce{Eu^2+}, cristalizou a partir dos fundidos que construíram a crosta superior e ficou
mais embaixo.
\end{solution}

\begin{solution}{exo:b2:elements-in-the-earth:5}
$\bar D_{\mathrm{La}} = 0.65\times0.058+0.18\times0.001+0.17\times0.001 = 0.0377 +
0.0002 + 0.0002 = 0.038$; $\bar D_{\mathrm{Lu}} = 0.3445+1.044+0.0306 = 1.42$. O lutécio é
compatível ($\bar D > 1$), o lantânio fortemente incompatível. A granada contém
$1.044/1.42 = 74$ por cento do lutécio da rocha, embora seja apenas 18 por cento da
rocha.
\end{solution}

\begin{solution}{exo:b2:elements-in-the-earth:6}
Para o sódio, a razão Na/Mg da Terra silicática é $2670/\num{228000} = 0.01171$ e o valor
CI, $5010/\num{95870} = 0.05226$: a razão é 0.224. Para o níquel, $1960/\num{228000}
= 0.008596$ contra $\num{10640}/\num{95870} = 0.1110$: 0.0775. O sódio é
empobrecido por sua volatilidade (\qty{958}{K}); o níquel se condensa acima de \qty{1300}{K} e é
três vezes mais empobrecido que o sódio porque é siderófilo: o núcleo levou a maior parte
dele.
\end{solution}

\begin{solution}{exo:b2:elements-in-the-earth:7}
La: $5.21\times0.058 = \qty{0.302}{ppm}$; Lu: $0.53\times0.53 = \qty{0.281}{ppm}$. Os
valores de normalização se cancelam na razão das razões:
$(\mathrm{La}/\mathrm{Lu})_{N,\mathrm{cpx}}/(\mathrm{La}/\mathrm{Lu})_{N,\mathrm{melt}} =
D_{\mathrm{La}}/D_{\mathrm{Lu}} = 0.058/0.53 = 0.11$: o padrão do cristal é nove
vezes mais inclinado em direção às terras raras pesadas que o do fundido.
\end{solution}

\begin{solution}{exo:b2:elements-in-the-earth:8}
Para a crosta continental superior, obtém-se $(12/0.658)/(31/0.648) = 18.24/47.84 = 0.38$. Basalto havaiano:
$(18.10/0.658)/(15.20/0.648) = 27.51/23.46 = 1.17$. MORB: $7.96/8.04 = 0.99$. Só a
crosta, com um forte déficit de nióbio, traz a assinatura das zonas de subducção.
\end{solution}

\begin{solution}{exo:b2:elements-in-the-earth:9}
\omterm{def:b2:elements-in-the-earth:bse}{Manto primitivo}: $240/0.0203 = \num{11800}$. MORB: K $= 0.160\times0.830\times10^4 =
\qty{1328}{ppm}$, K/U $= \num{11200}$. Crosta: K $= 2.80\times0.830\times10^4 =
\qty{23240}{ppm}$, K/U $= \num{8600}$. O potássio aumenta cem vezes do manto para a crosta,
e no entanto K/U varia menos de um terço: os dois elementos são quase igualmente incompatíveis,
e a fusão quase não os separa. O urânio é refratário, portanto sua abundância decorre
do modelo condrítico; a razão K/U medida em lavas dá então o potássio do planeta, que,
sendo volátil, não pode ser tirado dos \omterm{def:b2:planet-formation-meteorites:chondrule}{condritos}.
\end{solution}

\begin{solution}{exo:b2:elements-in-the-earth:10}
No equilíbrio, $C_{\mathrm{rock}} = \bar D\,C_{\mathrm{melt}}$ para cada elemento (e uma
pequena fração de fundido deixa a composição da rocha quase inalterada), portanto
$(C_1/C_2)_{\mathrm{melt}} = (C_{1}/\bar D_1)/(C_{2}/\bar D_2)$ para os $C_1,
C_2$ da rocha, e a razão é $\bar D_2/\bar D_1$. Para La e Lu: $1.419/0.03805 = 37$. O fundido
de um manto com granada é 37 vezes mais rico em La em relação ao Lu que sua fonte: a granada
reteve as terras raras pesadas, e o padrão do fundido é inclinado.
\end{solution}

\begin{solution}{exo:b2:elements-in-the-earth:11}
$C_{\mathrm{mix}} = fC_U + (1-f)C_M$ para cada elemento: linear em $f$. A razão
$\bigl(f\,\mathrm{Th}_U + (1-f)\mathrm{Th}_M\bigr)/\bigl(f\,\mathrm{Nb}_U +
(1-f)\mathrm{Nb}_M\bigr)$ é um quociente de duas funções lineares, uma hipérbole em $f$.
Th/Nb $= 0.2$ quando $f(10.5-2.4) + (1-f)(0.404-1.048) = 0$, isto é, $8.1f = 0.644(1-f)$,
$f = 0.644/8.744 = 0.074$: cerca de 7 por cento de crosta elevam Th/Nb de 0.077 para 0.2. Para
$f = 0.1$: Th $= 1.05 + 0.364 = 1.414$, Nb $= 1.2 + 4.716 = 5.916$, Th/Nb $= 0.24$.
\end{solution}

\begin{solution}{exo:b2:elements-in-the-earth:12}
Remover um componente que não contém nenhum dos dois elementos deixa as duas massas, e portanto sua razão,
inalteradas. Se todo o alumínio foi conservado, $(\mathrm{Al}/\mathrm{Mg}) = 1.16$ vezes o CI
significa que a Terra silicática contém $1/1.16 = 0.86$ de seu magnésio condrítico: faltam 14 por
cento. O silício em relação ao alumínio vale $0.83/1.16 = 0.72$ vez o valor CI:
faltariam 28 por cento do silício condrítico.
\end{solution}

\begin{solution}{pb:b2:elements-in-the-earth:1}
\textbf{1.} Fe/Mg $= \num{182800}/\num{95870} = 1.907$; Ni/Mg $= 0.1110$; Si/Mg $=
1.111$; Al/Mg $= 0.0887$.

\textbf{2.} $(\num{182800}/55.85)/(\num{106500}/28.09) = 3273/3791 = 0.863$ átomo de ferro por
átomo de silício, a menos de 7 por cento dos 0.81 da fotosfera: razoável, pois os \omterm{def:b2:elements-in-the-earth:ci}{condritos CI}
correspondem ao Sol nos elementos formadores de rochas.

\textbf{3.} A perda de um componente que não contém nenhum dos elementos de uma razão deixa a
razão inalterada (\cref{exo:b2:elements-in-the-earth:12}): só as concentrações absolutas mudam.

\textbf{4.} O magnésio é litófilo e não entra no núcleo; é um elemento maior,
bem determinado tanto nos \omterm{def:b2:planet-formation-meteorites:chondrule}{condritos} quanto no manto. O ferro é o elemento que se reparte
entre núcleo e manto: ele não pode ser sua própria referência.

\textbf{5.} Elementos que se condensam juntos não foram separados pela volatilidade na
nebulosa, portanto a Terra total deve conservar suas razões. Al/Mg se afasta: o alumínio é
refratário (\qty{1653}{K}), e a razão Al/Mg da Terra silicática vale 1.16 vezes o valor CI.

\textbf{6.} Fe $= 8.05\times55.85/71.85 = 6.26$ por cento; Ni $= 0.25\times58.69/74.69 =
0.196$ por cento.

\textbf{7.} Fe/Mg $= 6.26/22.8 = 0.2746$; Ni/Mg $= 0.196/22.8 = 0.00860$.

\textbf{8.} Ferro: $1 - 0.2746/1.907 = 0.856$, faltam 85.6 por cento; níquel:
$1 - 0.00860/0.1110 = 0.923$, 92.3 por cento.

\textbf{9.} O irídio é mais fortemente siderófilo que o níquel, e o níquel mais que o ferro:
quanto mais um elemento prefere o metal, mais dele o metal levou para o núcleo.

\textbf{10.} Nos silicatos, como \ce{Fe^2+} na olivina e no piroxênio (8 por cento de FeO):
o ferro oxidado a FeO liga-se ao oxigênio como um \omterm{def:b2:elements-in-the-earth:goldschmidt}{elemento litófilo} e permaneceu no
manto; só o ferro metálico afundou.

\textbf{11.} $(1-x)C_{\mathrm{Fe}} + xc_{\mathrm{Fe}} = r_{\mathrm{Fe}}(1-x)C_{\mathrm{Mg}}$;
dividir por $(1-x)c_{\mathrm{Fe}}$ dá a proposição.

\textbf{12.} $r_{\mathrm{Fe}}C_{\mathrm{Mg}} - C_{\mathrm{Fe}} = 1.907\times22.8 - 6.26 =
37.22$; $y = 37.22/85.5 = 0.4353$, $x = y/(1+y) = 0.303$.

\textbf{13.} $0.1110\times22.8 - 0.196 = 2.335$; $y = 2.335/5.2 = 0.4490$, $x = 0.310$.

\textbf{14.} $y = (85.5\times37.22 + 5.2\times2.335)/(85.5^2 + 5.2^2) =
3194.4/7337.3 = 0.4354$, $x = 0.303$.

\textbf{15.} Cada equação pesa como o quadrado de seu coeficiente: $85.5^2/5.2^2 = 270$
vezes mais para o ferro, o elemento muito mais abundante.

\textbf{16.} $37.22/2.335 = 15.9$, contra $85.5/5.2 = 16.4$ no modelo: uma diferença de 3 por
cento.

\textbf{17.} $0.303\times\num{5.972e24} = \qty{1.81e24}{kg}$.

\textbf{18.} $(32.5-30.3)/32.5 = 7$ por cento menor.

\textbf{19.} $y = 0.325/0.675 = 0.4815$, Fe no núcleo $= 37.22/0.4815 = 77.3$ por cento,
e Ni $= 2.335/0.4815 = 4.85$ por cento: 18 por cento de elementos leves, cerca do dobro dos 5 a
10 por cento permitidos pela densidade. A razão Fe/Mg da Terra deve ser um pouco mais alta que o valor
CI, ou seu magnésio silicático mais alto que o do modelo.

\textbf{20.} Silício, enxofre, carbono, oxigênio e hidrogênio. O silício está em falta na
Terra silicática: sua razão normalizada ao Mg, 0.83, fica abaixo do nível refratário.

\textbf{21.} $1.111\times22.8 - 21.0 = 4.331$; $4.331/0.4354 = 9.9$ por cento de Si no
núcleo.

\textbf{22.} 9.9 contra 6 por cento: o núcleo pode conter mais silício que o modelo, ou
parte do silício nunca foi acretada. O silício se condensa a \qty{1310}{K}, um pouco
abaixo do magnésio e do ferro (1336 e \qty{1334}{K}): parte dele pode ter ficado no gás.

\textbf{23.} Núcleo: $0.325\times0.5 = \qty{0.1625}{ppm}$ da massa da Terra; Terra
silicática: $0.675\times0.001 = \qty{0.000675}{ppm}$. O núcleo contém
$0.1625/0.1632 = 99.6$ por cento do ouro.

\textbf{24.} A química sozinha coloca
\textbf{\qty[text-series-to-math=true]{30.3}{\percent}} da massa da Terra em seu núcleo,
a menos de 7 por cento dos 32.5 por cento da sismologia.
\end{solution}
